\begin{frame}{제약 문제와 투영의 그림}
  \begin{columns}[T]
    \begin{column}{0.52\textwidth}
      \begin{block}{PCA 문제를 하나의 문장으로}
        \[ f(v) = v^{T}\Sigma v \;\text{를 최대화},\;
           \text{제약 } v^{T}v = 1 \]
        $v$를 키우면 $v^{T}\Sigma v$도 같이 커지므로,
        \textbf{방향만} 비교하려면 $\|v\|=1$ 제약을 건다.
      \end{block}
      \vskip0.6em
      \begin{itemize}
        \item 제약 최댓값 문제의 표준 도구: \kb{라그랑주 승수법}.
      \end{itemize}
    \end{column}
    \begin{column}{0.46\textwidth}
      \centering
      \includegraphics[width=\linewidth]{ch14_pca_projection.png}
      \vskip0.3em
      {\footnotesize 4개 점을 분산이 최대인 PC1(45도 대각선, 분산 4.5)으로
      투영. x축(분산 2.5)보다 대각선 방향이 정보를 훨씬 많이 담는다.}
    \end{column}
  \end{columns}
\end{frame}
